% !PRINTED 377
\weberpage{377}
%<<CONT
durch den Schnitt von $x_1$ und $x_4$ gehen, und folglich ist $x_4$ aus $x_1$ und $\lambda_2 x_3 + \dfrac{x_2}{\lambda_3}$ linear zusammengesetzt. Nun aber können wir $x_4$ linear durch $x_1, x_2, x_3$ ausdrücken, etwa in der Form

\weq{10}{x_4 = a_1 x_1 + a_2 x_2 + a_3 x_3,}

und darin sind $a_1, a_2, a_3$ als bekannt zu betrachten, und keine dieser Constanten kann verschwinden. Es ist also $\lambda_2\lambda_3 = a_3 : a_2$, und wenn wir eine neue Constante $h_1$ einführen, so ist

\weq{11}{\lambda_2 = h_1 a_3,\quad \frac{1}{\lambda_3} = h_1 a_2,}

und die Gleichung (9) ergiebt die Identität:

\[
h_1 x_4\left(\frac{\xi_{42}}{\lambda_2} + \frac{\xi_{43}}{\lambda_3} - h_1\xi_1\right) = x_1\left(-h_1(2+a_1h_1)\xi_1 + \frac{\xi_2}{a_2} + \frac{\xi_3}{a_3}\right).
\]

Daraus schliesst man weiter, wenn $k_1$ eine neue Constante ist,

\[
k_1 x_4 = -h_1(2+a_1h_1)\xi_1 + \frac{\xi_2}{a_2} + \frac{\xi_3}{a_3},
\]

\weq{12}{\frac{\xi_{42}}{\lambda_2} + \frac{\xi_{43}}{\lambda_3} = h_1\xi_1 + \frac{k_1}{h_1} x_1.}

Wenn man hierin die Ziffern 1, 2, 3 cyklisch vertauscht, so ergeben sich aus (8) drei solche Systeme; zunächst:

\weq{13}{\begin{aligned}
k_1 x_4 &= -h_1(2+a_1h_1)\xi_1 + \frac{\xi_2}{a_2} + \frac{\xi_3}{a_3}\\
k_2 x_4 &= \frac{\xi_1}{a_1} - h_2(2+a_2h_2)\xi_2 + \frac{\xi_3}{a_3}\\
k_3 x_4 &= \frac{\xi_1}{a_1} + \frac{\xi_2}{a_2} - h_3(2+a_3h_3)\xi_3.
\end{aligned}}

Da diese drei Ausdrücke für $x_4$ mit einander identisch sein müssen, so folgt:

\[
-\frac{h_1(2+a_1h_1)}{k_1} = \frac{1}{k_2 a_1} = \frac{1}{k_3 a_1},
\]

also $k_2 = k_3$ und ebenso $k_2 = k_1$. Es sind also $k_1, k_2, k_3$ einander gleich und wir setzen dafür $k$. Dann folgt weiter

\[
a_1 h_1(2+a_1h_1) + 1 = (a_1h_1+1)^2 = 0,
\]

also $h_1 a_1 = -1$, und ebenso

\[
h_1 = \frac{-1}{a_1},\quad h_2 = \frac{-1}{a_2},\quad h_3 = \frac{-1}{a_3}.
\]
